\subsection*{All abelian refinements of the actual product}

\statement{B.58}{Proved}{Rigid-seven tests detect ambient abelianization integrally}
Let $P=\Gamma_{0,12}$, and let $\mathscr S_7$ be the twelve
rigid seven-puncture subsets in B.32 and B.56.
The combined forgetful homomorphism induces a split injection
of free abelian groups
\[
 H_1(P,\mathbb Z)\cong\mathbb Z^{54}
 \longrightarrow
 \prod_{S\in\mathscr S_7}H_1(P\Gamma_{0,S},\mathbb Z)
 \cong\mathbb Z^{168}.                                      \tag{B.58.1}
\]
Tensoring with $\mathbb Z/m$ retains injectivity for every
$m\geq1$.  The same holds with $\widehat{\mathbb Z}$.
Consequently the actual residual centralizer satisfies
\[
 \mathcal K\subseteq
 \ker\bigl(\widehat P\longrightarrow
                    \widehat{P^{\mathrm{ab}}}\bigr)
 =\overline{[P,P]}.                                         \tag{B.58.2}
\]
For any two finite-index subgroups $L_1,L_2\leq P$ with
$[P,P]\subseteq L_1\cap L_2$, the product
\[
       (H_T\cap L_1)(H_5\cap L_2)                           \tag{B.58.3}
\]
is closed in the full profinite topology of $P$.
In particular all one-factor refinements
$(H_T\cap L)H_5$ with $[P,P]\subseteq L\leq_f P$ are closed.
This proves a family of finite-index refinements of B.52,
but does not prove the refinements associated to arbitrary
nonabelian finite quotients.

\begin{proof}
\emph{The integral lattice and its forgetful maps.}
Normalize three labelled points $a,b,c$ to
$\infty,0,1$.  Then $M_{0,n}$ is the configuration space
of $n-3$ labelled points of $\mathbb C\setminus\{0,1\}$.
Its integral first homology is freely detected by the
winding numbers of each variable point about $0$ and $1$,
and by each difference of two variable coordinates.
One direct proof uses the forgetful fibrations: their fibres
are punctured planes, their pure monodromies act trivially
on the first homology of the fibre, and they have sections
obtained by adding a point sufficiently far away.
Induction gives precisely
$2(n-3)+\binom{n-3}{2}=n(n-3)/2$ independent coordinates.
Equivalently these are the pair linkings of the $n-1$
finite normalized coordinates, omitting the constant pair
$(b,c)$.  These coordinates describe the pure sphere
\emph{mapping class} group; no central order-two sphere
braid factor is retained.

Write $e_{ij}=e_{ji}$ for a formal original pair coordinate.
For distinct finite normalized labels $i,j$, the corresponding
linking covector in the original coordinates is
\[
 v(a,b,c;i,j)
   =e_{ca}-e_{cb}+e_{ba}+e_{ij}-e_{ia}-e_{ja}.                \tag{B.58.4}
\]
Indeed, if the original coordinates are $x_i$, their
normalized difference, up to a constant sign, is
\[
 z_i-z_j=
 \frac{(x_c-x_a)(x_b-x_a)(x_i-x_j)}
      {(x_c-x_b)(x_i-x_a)(x_j-x_a)}.
\]
Taking winding numbers proves (B.58.4).
Cancellation also covers pairs incident to $b$ or $c$.
The omitted pair $(b,c)$ has covector zero.

For the ambient integral coordinates take $(a,b,c)=(2,4,6)$.
Let $\mathcal E$ be the $54$ unordered pairs of
$\{1,\ldots,12\}\setminus\{2\}$ other than $(4,6)$,
in lexicographic order.  Form the $54\times66$ matrix $A$
whose rows are $v(2,4,6;i,j)$ for $(i,j)\in\mathcal E$.
On the columns $\mathcal E$ its restriction is the identity.
For each $S\in\mathscr S_7$, in lexicographic order, choose
the three smallest labels as $a,b,c$ and form its fourteen
rows (B.58.4), also in lexicographic pair order.  This gives
a completely specified $168\times66$ matrix $R$.
Every row annihilates the sphere vertex relations.
If $M$ is the restriction of $R$ to its columns
$\mathcal E$, then direct integer arithmetic gives
\[
                         MA=R.                             \tag{B.58.5}
\]
Thus $M$ is the matrix of the actual homomorphism (B.58.1)
in the specified integral coordinates.

The exact certificate
\nolinkurl{check_ambient_rigid7_abelian.py}
constructs these matrices from the twelve subsets already
certified in \nolinkurl{rigid7_subsets.json}.
It verifies (B.58.5), computes rank $54$, and gives
$54$ Smith invariant factors all equal to $1$.
More concretely, the first $54$ independent rows of $M$
form a square matrix of determinant $-1$.
The complete matrix, the selected row indices, and its
\emph{integer inverse} are supplied in
\nolinkurl{ambient_rigid7_abelian_certificate.json}.

The supplied integer inverse is checked by ordinary
integer multiplication.  Composing it with the selected
coordinate projection gives an integral left inverse to $M$.
The independent program
\nolinkurl{code/verify_B58_certificate.py}
also reconstructs the subsets from the six order-five
subgroups and verifies both products of the displayed
inverse matrices; its standard output is retained.
The first program's actual standard output is
\nolinkurl{check_ambient_rigid7_abelian.out}.
The integral left inverse remains a left inverse over
every coefficient ring asserted in the statement.

\emph{The residual group and the refinements.}
B.32 shows that every element of $\mathcal K$ is killed
by each of these twelve rigid forgetful maps.
Passing to abelianizations and using the
$\widehat{\mathbb Z}$ version of (B.58.1) proves
(B.58.2).  The equality of its kernel with
$\overline{[P,P]}$ follows directly from the universal
property of the full profinite completion of the quotient
$P/[P,P]$.

Let $g\in P$ lie in the closure of (B.58.3).
B.52 first gives its unique discrete factorization
$g=hj$ with $h\in H_T$ and $j\in H_5$.
Compactness of the completed factors also gives
$g=\hat h\hat j$, with
$\hat h\in\overline{H_T\cap L_1}$ and
$\hat j\in\overline{H_5\cap L_2}$.
Then
\[
 h^{-1}\hat h=j\hat j^{-1}
 \in\overline H_T\cap\overline H_5=\mathcal K.
\]
By (B.58.2) this element belongs to both
$\overline L_1$ and $\overline L_2$.
Hence $h\in P\cap\overline L_1=L_1$ and
$j\in P\cap\overline L_2=L_2$.
Therefore $g$ belongs to (B.58.3), which proves its closure.

\emph{Tetrahedral cross-check.}
The separate certificate
\nolinkurl{check_full_pair_linking.py} computes the orders
of all $55$ rational differences of the eleven finite
normalized orbit points at all thirteen basis meridians
of $S_T$.  The calculations use the exact number field
$\mathbb Q(i,\sqrt3)$.
The restriction to the rigid seven labels agrees with
all thirteen independently certified braid-linking vectors
from B.56.  Both the full and restricted matrices have rank
six; an explicit integer matrix expresses six times the
full matrix in terms of six rows of the restricted matrix.
This proves the same abelian vanishing on the completed
tetrahedral subgroup, including at the primes two and
three, since multiplication by six on
$\widehat{\mathbb Z}$ is injective.
The complete rational differences, valuation matrices and
reconstruction are in
\nolinkurl{full_pair_linking_certificate.json} and the
actual output is \nolinkurl{check_full_pair_linking.out}.
This second computation is not needed for the stronger
ambient split-injection argument.
\end{proof}

In particular, none of the nonzero degree-one fixed classes
in $\mathbb Z_\ell^2$ from B.20 is the image of an element
of the actual $\mathcal K$.
This makes the failure of a degree-one fixed-point lifting
inference concrete.  It does not show that $\mathcal K$ is
trivial or perfect, and it does not determine its image
in every finite nonabelian quotient of $P$.
